% Generated draft: requires mathematical and source review.
\documentclass[11pt]{article}
\usepackage[margin=25mm]{geometry}
\usepackage{fontspec}
\setmainfont{Latin Modern Roman}
\usepackage{amsmath,amssymb,mathtools,unicode-math}
\setmathfont{Latin Modern Math}
\usepackage{enumitem,needspace,xurl,hyperref}
\setlength{\parindent}{0pt}
\setlength{\parskip}{5pt}
\setlength{\emergencystretch}{3em}
\newcommand{\N}{\mathbb{N}}
\newcommand{\Z}{\mathbb{Z}}
\newcommand{\Q}{\mathbb{Q}}
\newcommand{\R}{\mathbb{R}}
\newcommand{\C}{\mathbb{C}}
\newcommand{\F}{\mathbb{F}}
\newcommand{\A}{\mathbb{A}}
\newcommand{\PP}{\mathbb{P}}
\newcommand{\E}{\mathbb{E}}
\newcommand{\eps}{\varepsilon}
\newcommand{\dd}{\,\mathrm{d}}
\begin{document}
\section*{UnsolvedMath ID: 1361}
\textbf{Source problem code:} catalogue entry
\subsection{Definitions and mathematical statement}
The source asks, for a board of given width, for the least dimension in which the first player X has a winning strategy. It does not specify the game rules sufficiently to identify one numerical function.

A conditional formulation of the standard finite placement model is as follows. For integers $n,d\geq1$, let the board be $B_{n,d}=[n]^d$, where $[n]=\{1,\ldots,n\}$, and specify a family $\mathcal L_{n,d}$ of winning subsets of the board. Two players alternately claim one previously unclaimed cell, with X moving first. A player wins upon claiming every cell of a member of $\mathcal L_{n,d}$; if the board is filled without a winner, the game is drawn. A winning strategy for X is a rule for legal moves, depending on the preceding moves, that forces X to win against every legal strategy of the second player. For these specified rules, the intended quantity would be
\[
d_*(n)=\min\{d\geq1:\text{X has a winning strategy on }(B_{n,d},\mathcal L_{n,d})\},
\]
with $d_*(n)=\infty$ if the set is empty.

The choice of $\mathcal L_{n,d}$ is material. For example, a combinatorial line has one nonempty set of variable coordinates all equal to a common $t\in[n]$, with the other coordinates fixed. A model admitting reversed coordinates $n+1-t$ as well has a different family of lines. Further variants ask for a fixed number of consecutive cells on a wider board. These alternatives are not selected by the stored wording. The board $[n]^d$ and alternating-placement rules above are therefore an explicit conditional model, not asserted missing source data.

The background says ordinary tic-tac-toe is a draw with perfect play and invokes the Hales--Jewett theorem for existence, at fixed line length $n$, of a sufficiently large dimension in which the first player wins. That explanation depends on adopting the relevant line family and standard symmetric placement rules; it does not specify an exact threshold for the undefined stored game. The supplied triage dated 2026-08-17 says that geometry, winning length, and whether width remains fixed across dimensions are missing and that several distinct multidimensional games fit the sentence. It lists no separately verified partial result, despite the partially-solved triage label; the registry status is open. No optimum or general threshold is asserted here.
\subsection{Short English statement}
For each board width, determine the least dimension in which X can force a win. A precise question requires the board, winning lines or winning length, and move rules; the supplied wording leaves these unresolved, so it does not define a unique threshold.
\subsection{Sources}
\begin{itemize}
\item ``Tic-Tac-Toe Winning Dimension,'' source entry cited by the supplied triage, accessed 2026-08-17 according to the record: \url{https://www.unsolvedmath.com/}. No original Hales--Jewett citation or separate authoritative game-rules reference is supplied.
\item Local registry data/unsolvedmath/problems.json, permanent catalogue entry, source code catalogue entry, ``Tic-Tac-Toe Winning Dimension''; statement, background, and literature triage dated 2026-08-17. Dataset attribution: Ulam AI and the UnsolvedMath Contributors.
\end{itemize}
\end{document}
